\section{Setting and main results} \label{sec:setting-results}

In the following, Cartesian products are always equipped with product $\sigma$-algebras. For finite interacting particle systems on lattices, the configuration space is $\N_0^\Lambda = \set{(n_x)_{x\in \Lambda}\colon n_x \in \N_0}$ with $\Lambda\subset\Z^d$. We replace $\N_0^\Lambda$ with a space of finite counting measures:
Let $(E,\mathcal E)$ be a Borel space, for example, $E= \R^d$ or more generally a Polish space with its Borel $\sigma$-algebra. Let~$\mathbf N_{<\infty}$ be the space of finite counting measures on $E$. Every element $\eta \in \mathbf N_{<\infty}$ is either zero or a finite sum $\eta= \delta_{x_1}+\cdots + \delta_{x_n}$ of Dirac measures, where $x_1, \ldots, x_n \in E$, $n \in \N$. The space $\mathbf N_{<\infty}$ is equipped with the $\sigma$-algebra $\mathcal N_{<\infty}$ generated by the counting variables $\eta \mapsto \eta(B)$, $B\in \mathcal E$.
For background on point processes and counting measures, we refer the reader to Last and Penrose~\cite{LastPenroseLecturesOnThePoissonProcess}.

\subsection{Symmetric inclusion process}

We are interested in Markov processes with state space $\mathbf N_{<\infty}$. As a guiding example we take the process with formal generator
\begin{equation} \label{eq:gsip}
	L f(\eta) = \iint \bigl( f(\eta-\delta_x+\delta_y)- f(\eta)\bigr) c(x,y) \eta(\dd x) \bigl( \alpha+\eta\bigr) (\dd y),
\end{equation}
where $\alpha$ is a finite measure on $E$ and $c\colon E\times E\to \R_+$ is a bounded measurable function with $c(x,y) = c(y,x)$ on $E^2$. In \cite{IntertwiningConsistent} we called this process \emph{generalized inclusion process}. When $E$ is a~finite set, we may identify finite counting measures $\eta$ with vectors $\vect n = (n_x)_{x\in \Lambda}\in \N_0^\Lambda$, integrals over $E$ turn into sums, and the generator turns into
\begin{equation} \label{eq:sip}
	Lf(\vect n) = \sum_{x,y\in E} \bigl( f(\vect n- \vect e_x+\vect e_y) - f(\vect n)\bigr) c(x,y) n_x (\alpha_y+n_y),
\end{equation}
where $\vect e_x(y) = \delta_{x,y}$ and $\alpha_y = \alpha(\{y\})$. This is the generator of the inhomogeneous \emph{symmetric inclusion process} (see, e.g., \cite[equation~(2.2)]{floreani2020orthogonal}), which first appeared in the homogeneous case as a dual process to a model for energy and momentum transport, see \cite{giardina-redig-vafayi2010} and references therein.

When $c(x,y)$ is the constant function with value $1$, the process corresponds to the Moran process from mathematical population genetics--classical Moran model \cite{TransitionMoran,Moran} when $E$ is finite, measure-valued Moran model (see \cite[Section 2.6]{DawsonMeasureValuedMarkovProcesses} or \cite[Section~5.4]{Etheridge00}) when $E$ is uncountable. The Moran model describes a population in which individuals have genetic types $x\in E$ that may mutate or evolve by sampling replacement (death of an individual followed by replacement with the offspring of another individual).

For the homogeneous symmetric inclusion process with finite state space / Moran model with finite type space, the correspondence was already noticed in \cite[Section~5]{carinci2015dualities}. For uncountable state spaces, we notice that the $n$-particle dynamics of the process generated by \eqref{eq:gsip} (with $c(x,y) \equiv 1$) coincides with the $n$-particle Moran model \cite[equation~(2.5.2)]{DawsonMeasureValuedMarkovProcesses} with so-called \emph{mutation operator} $A\varphi(x) = \int( \varphi(y) - \varphi(x)) \alpha(\mathrm d x)$. Our process for counting measures is similar to the measure-valued Moran model. The only difference is in bookkeeping: In mathematical population genetics, the population is often modeled not with $\eta \in \mathbf N_{<\infty}$ but rather with the empirical measure $\frac{\eta}{\eta(E)}$. This is helpful for scaling limits in which the population size goes to infinity, notably Fleming--Viot limits, which are important in population genetics (see \cite{shigaDiffusions} and references therein).

\subsection{Pascal point process} \label{sec:pascaldef}
The symmetric inclusion process with generator~\eqref{eq:gsip} has a family of reversible measures $\rho_{p,\alpha}$ indexed by $p\in (0,1)$ \cite[Theorem 5.2]{IntertwiningConsistent} where $\rho_{p,\alpha}$ is the law of the Pascal point process or negative binomial process (see, e.g., \cite[Section~2.7]{serfozo1990point}, \cite{kozubowski2008distributional} and \cite[Section~5.2]{IntertwiningConsistent}). We recall the definition here.
For $a\in \R$ and $n\in \N$, the rising factorial (also known as Pochhammer symbol)~is
\[
	(a)_0 = 1,\qquad (a)_n = a(a+1)\cdots (a+n-1).
\]

\begin{Definition}
 Let $p\in (0,1)$ and $\alpha$ be a finite measure on $E$. $\rho_{p,\alpha}$ is the uniquely defined probability measure on $\mathbf N_{<\infty}$ such that:
 \begin{itemize}\itemsep=0pt
 	\item For all $B\in \mathcal E$, and $n\in \N_0$
 	\begin{equation} \label{eq:negbin}
 		\rho_{p,\alpha}(\{\eta\colon \eta(B) =n\}) = (1-p)^{\alpha(B)} \frac{p^n}{n!} (\alpha(B))_{n}.
 	\end{equation}
 	\item For all $n \geq 2$ and all disjoint $B_1,\ldots, B_n\in \mathcal E$, the variables $\eta\mapsto \eta(B_i)$, $i=1,\ldots, n$, are independent.
 \end{itemize}
\end{Definition}

We work in the Hilbert space $\mathcal H = L^2(\mathbf N_{<\infty}, \mathcal N_{<\infty}, \rho_{p,\alpha})$ of complex-valued square-integrable functions. The scalar product is
\[
	\langle f,g \rangle = \int \overline{f} g\, \dd\rho_{p,\alpha}.
\]

\subsection[Representation of the su(1,1)]{Representation of the $\boldsymbol{\mathfrak{su}(1,1)}$ current algebra} \label{sec:su11}

Let $\mathcal D\subset \mathcal H$ be the set of bounded measurable functions $F$ for which there exists a cutoff $ m_F \in \N$ such that $F$ is supported in $\{\eta\colon \eta(E)\leq m_F\}$. We define operators $k^\pm(\varphi),k^0(\varphi)\colon\mathcal D\to \mathcal D$ indexed by bounded measurable functions $\varphi\colon E\to \C$ as follows:
\begin{align*}
	&k^+(\varphi) F(\eta) = \frac{1}{\sqrt p} \int \varphi(x) F(\eta-\delta_x) \eta(\dd x),\\
 &k^-(\varphi) F(\eta) = \sqrt{p} \int \overline{\varphi(x)} F(\eta+\delta_x) \bigl(\alpha+\eta)\bigr(\dd x),\\
	&k^0(\varphi) F(\eta) = F(\eta) \biggl( \int \varphi\,\dd \eta+ \frac12 \int \varphi \,\dd \alpha \biggr).
\end{align*}
for $\eta \in \mathbf{N}_{<\infty}$.
The raising operator $k^+(\varphi)$ and the lowering operator $k^-(\varphi)$ map functions supported in $\{\eta\colon \eta(E) = n\}$ to functions supported in $\{\eta\colon \eta(E) =n+1\}$ and $\{\eta\colon \eta(E) =n-1\}$ respectively. The indicator $\1_{\{0\}}$ that there is no particle at all (\emph{vacuum}) is annihilated by all lowering operators: $k^-(\varphi) \1_{\{0\}}=0$.

Our operators are a representation of the $\mathfrak{su}(1,1)$ current algebra, i.e., they represent the matrix-valued maps $x\mapsto \varphi(x) k^+$, $x\mapsto \overline{\varphi(x)} k^-$, $x\mapsto \theta(x) k^0$ with
\begin{align*}
 	k^+ = \begin{pmatrix} 0 & \mathrm i \\ 0 & 0 \end{pmatrix},\qquad
	k^- = \begin{pmatrix} 0 & 0 \\ \mathrm i & 0 \end{pmatrix},\qquad
	k^0 = \begin{pmatrix} 1/2 & \hphantom{-}0 \\ 0 & -1/2 \end{pmatrix},
\end{align*}
see \cite[Section~2.2]{algebraic_Floreani}.
The following commutation relations can be checked by an explicit computation:
\begin{gather*}
	\bigl[k^-(\varphi), k^+(\theta)\bigr] = 2 k^0(\overline{\varphi}\theta),\qquad
\bigl[k^0(\theta),k^+(\varphi)\bigr] = k^+(\varphi\theta),\qquad [k^0(\theta),k^-(\varphi)] = - k^-(\overline{\varphi}\theta),
\end{gather*}
where $[T, S] := TS - ST$.
Indeed, it can be readily verified that
\begin{gather*}
 k^-(\varphi) k^+(\theta) F(\eta) = \iint \overline{\varphi(y)} \theta(x) F(\eta - \delta_x + \delta_y) \eta(\dd x) (\alpha + \eta)(\dd y) \\
 \hphantom{k^-(\varphi) k^+(\theta) F(\eta) = \iint}{}
 + F(\eta) \int \overline{\varphi(y)} \theta(y) (\alpha + \eta)(\dd y), \\
 k^+(\theta) k^-(\varphi) F(\eta) = \iint \overline{\varphi(y)} \theta(x) F(\eta - \delta_x + \delta_y) \eta(\dd x) (\alpha + \eta)(\dd y) \\
 \hphantom{k^+(\theta) k^-(\varphi) F(\eta) =\iint}{}
 + F(\eta) \int \overline{\varphi(x)} \theta(x) \eta(\dd y).
\end{gather*}
Subtracting these equations yields the commutation relation
\[
\bigl[k^-(\varphi), k^+(\theta)\bigr] F(\eta) = F(\eta) \int \overline{\varphi(x)} \theta(x) (\alpha + 2\eta)(\dd x) = 2k^0(\overline{\varphi}\theta) F(\eta).
\]
 The other relations follow similarly.

The factor $\sqrt p$ included in the definitions of $k^\pm(\varphi)$ is irrelevant for the commutation relations but it matters for the adjointness relation $\la f, k^+(\varphi) g\ra = \la k^-(\varphi) f,g\ra$ proven in Lemma~\ref{lem:adjoints}.

We leave a thorough analysis of the current algebra generated by the operators $k^\pm(\varphi)$, $k^0(\varphi)$ to another article \cite{algebraic_Floreani} and focus on the operators associated with the constant function $\varphi = \1$, equal to $1$ on all of $E$. In Lemma~\ref{lem:closure} below, we check that for every $\xi \in \C$, the operator $\frac{1}{\mathrm i} \bigl( \xi k^+(\1) - \overline{\xi} k^-(\1))$ with domain $\mathcal D$ is closable with self-adjoint closure $A$. The operator $\exp( \xi k^+(\1) - \overline{\xi} k^-(\1))$ is, by definition, the unitary $\exp(\mathrm i A)$.
For $\theta \in \R$, the unitary $\exp\bigl( 2\mathrm i \theta k^0(\1)\bigr)$ is defined as a multiplication operator, it multiplies a function $f(\eta)$ with $\exp(\mathrm i \theta (\alpha (E) + 2 \eta(E)))$.
In this way we obtain a family of unitary operators
\begin{equation}\label{eq:unitary}
	{\rm U}(\xi,\theta) = \exp\bigl( \xi k^+(\1) - \overline{\xi} k^-(\1)\bigr) \exp\bigl(2\mathrm i\theta k^0(\1)\bigr),\qquad \xi \in \C,\theta\in \R.
\end{equation}
For the definition of the unitary~\eqref{eq:unitary}, it is essential that $\alpha$ has finite total mass and that each configuration $\eta$ has finitely many particles.

\subsection{Symmetries of consistent particle processes}

Let $(P_t)_{t\geq 0}$ be the semigroup of a Markov process with state space $\mathbf N_{<\infty}$. For a measurable $f\colon \mathbf N_{<\infty} \to \R_+$, let $\mathcal A f\colon \mathbf N_{<\infty} \to \R_+$ be the function given by
\[
	\mathcal A f(\eta) = \int f(\eta-\delta_x)\eta(\dd x).
\]

\begin{Definition}
 The semigroup $(P_t)_{t\geq 0}$ is \emph{consistent} if $P_t \mathcal A f = \mathcal A P_t f$ for all measurable $f\colon \mathbf N_{<\infty}\to \R_+$ and all $t\geq 0$.
\end{Definition}

Notice that, up to questions of domains, $\mathcal A = \sqrt p\, k^+(\1)$. In Proposition~\ref{proposition: consistency implies conservative} below, we show that if a process is consistent then it is \emph{conservative}, i.e., the total number of particles is constant in time. In view of that, we could normalize $\mathcal A$ by the total mass $\eta(E)$ (which is constant in time) and read consistency as the property that the action of removing a particle uniformly at random commutes with the dynamics \cite{carinci2021consistent,IntertwiningConsistent}.

The following theorem generalizes item 1\,(i) in Theorem~3.1 in Carinci et al.~\cite{carinci2019orthogonal}.

\begin{Theorem} \label{thm:symmetry}
Let $(P_t)_{t\geq 0}$ be the semigroup of a Markov process with state space $\mathbf N_{<\infty}$. Let $p\in (0,1)$ and $\alpha$ be a finite measure on $E$. Assume that the process is consistent, and admits the Pascal law $\rho_{p,\alpha}$ as a reversible measure. Then $P_t$ commutes with all unitaries ${\rm U}(\xi,\theta)$ from~\eqref{eq:unitary}:
	\[
		P_t {\rm U}(\xi,\theta) f(\eta) = {\rm U}(\xi,\theta) P_t f (\eta)
	\]
	for all $t\geq 0$, $f\in L^2(\mathbf N_{<\infty},\mathcal N_{<\infty},\rho_{p,\alpha})$ and $\rho_{p,\alpha}$-almost all $\eta\in \mathbf N_{<\infty}$.
\end{Theorem}

The generalized symmetric inclusion process is a consistent Markov process and admits the Pascal law $\rho_{p,\alpha}$ as a reversible measure (as proved in \cite[Theorem~5.2.]{IntertwiningConsistent}) and, thus, Theorem~\ref{thm:symmetry} applies.
